\begingroup
\makeatletter\tagsleft@true\let\veqno\@@leqno\makeatother
\fancyhead{}\fancyhead[R]{\small261}\fancyhead[C]{\scriptsize BEMERKUNGEN ZUM AGM. ABSCHNITT II.}
\SourceLeaf{268b}
\pdfbookmark[1]{Erläuterungen zu [II.]}{x1-agm-erlaeuterung-ii}
\begin{center}\textls[120]{Erläuterungen zu [II.], S. 177—180.}\end{center}
\par\indent Die Formeln [1], [5], [17] geben die bekannte\textsuperscript{***)} Reihenentwicklung des Ellipsenquadranten $q$ mit den Halbachsen $a,b$ $(a>b)$ nach Potenzen der Größe
\[B=\sqrt{\frac{a^2-b^2}{a^2}}\]
\noindent (siehe \textit{Abriß}, Gl. (21), S. 254, wo $k=B$ zu nehmen ist).
\par\indent In den Formeln [1], [5] ist der Faktor $\frac\pi2a$ unterdrückt. Die Gleichung [18] stellt, wenn man beiderseits den Faktor $a$ wegläßt, die Entwicklung von $\frac{a}{M(a,b)}$ nach Potenzen von $B$ dar (vergl. \textit{Abriß}, Gl. (19), S. 254). \textsc{Gauss} dürfte diese Entwicklung in derselben Weise gefunden haben, wie er sie im art. 7 der Abhandlung von 1800 (Werke III, S. 361, siehe dort S. 367) abgeleitet hat. Wir geben diese Ableitung hier wieder und verweisen auf eine zweite Art der Herleitung, die in der Aufzeichnung [III.] enthalten ist, und auf die wir weiter unten zurückkommen.
\par\indent Es ist
\[\frac{a}{M(a,b)}=\frac{a}{M(a+c,a-c)}=\frac1{M(1+B,1-B)}=\frac1{M(1,\sqrt{1-B^2})};\]
\noindent setzt man nun
\[B=\frac{2t}{1+t^2},\]
\vfill
\par\smallskip\noindent\rule{.25\linewidth}{.4pt}\par
{\small\par\indent ***) Man vergl. die oben S. 254 Fußnote angeführte Abhandlung \textsc{Eulers}, wo S. 128 die Reihenentwicklung für den Ellipsenquadranten genau in der von \textsc{Gauss} benutzten Form gegeben wird. Die Übereinstimmung auch in den Bezeichnungen bestätigt, daß \textsc{Gauss} jene Abhandlung \textsc{Eulers} gekannt hat; \textsc{Euler} sowohl wie \textsc{Gauss} haben $q$, statt des \textsc{Euler}schen $n$ hat \textsc{Gauss} in den Formeln des art. [1.] $B$, in späteren Formeln der artt. [2.] und [5.] aber $\nu$; den Buchstaben $n$ dürfte \textsc{Gauss} vermieden haben, weil er mit $n$ eine ganze Zahl zu bezeichnen pflegte. Auf den ersten Durchschußblättern des \textsc{Leiste} hat sich \textsc{Gauss} (wie es scheint nach der \textsc{Fuss}schen Liste vom Jahre 1783) ein Verzeichnis der Werke \textsc{Eulers} angelegt, er hat sich also schon sehr früh lebhaft für \textsc{Euler} interessiert.\par}
\clearpage
\ArithmeticHead{262}{L}{ANALYSIS. NACHLASS.}
\SourceLeaf{269}
\noindent so wird
\[M(1+B,1-B)=M(1,\sqrt{1-B^2})=M\left(1,\frac{1-t^2}{1+t^2}\right)=\frac1{1+t^2}M(1+t^2,1-t^2).\]
\noindent Setzt man ferner in der Funktionalgleichung
\[M\left(1+\frac{2t}{1+t^2},\,1-\frac{2t}{1+t^2}\right)=\frac1{1+t^2}M(1+t^2,1-t^2)\]
\noindent die Entwicklung
\[\frac1{M(1+B,1-B)}=1+\alpha B^2+\beta B^4+\gamma B^6+\ldots\]
\noindent mit unbestimmten Koeffizienten ein, so findet man
\[\alpha=\frac14,\ \beta=\frac14\,.\,\frac9{16},\ \gamma=\frac14\,.\,\frac9{16}\,.\,\frac{25}{36},\ldots\]
\par\indent Den allgemeinen Zusammenhang zwischen dem agM. und dem vollständigen elliptischen Integral erster Gattung hat \textsc{Gauss} zu der Zeit, wo er jene Formeln in den \textsc{Leiste} einschrieb, jedenfalls noch nicht gekannt (siehe weiter unten).
\par\indent Die linke Seite der Formel [2] und das erste Glied von [6] geben also die Entwicklung von
\[\frac{Ba}{M(a,b)}=\frac{\sqrt{a^2-b^2}}{M(a,b)}=\frac{c}{M(a,b)}.\]
\par\indent Die Gleichung [6] entspricht somit der Formel (33) des \textit{Abrisses}, wenn man dort $x=z^2$ setzt. Sie enthält also den Zusammenhang zwischen dem agM. und den Reihen, deren Exponenten eine arithmetische Reihe zweiter Ordnung bilden, bezw. die Quadratzahlen sind\textsuperscript{*)}.
\par\indent Die Gleichung [2] enthält auf der rechten Seite nur die zwei ersten Glieder der Reihe $r^2$, die Gleichung [3] entsteht aus [2] bezw. [6], indem man $B$ ausrechnet; in solchen Operationen mit Reihen, deren Inversion und dergl. war \textsc{Gauss} schon in sehr früher Zeit wohlgeübt\textsuperscript{**)}.
\par\indent Übrigens ist der Koeffizient von $z^5$ in [3] unrichtig, die Reihe lautet
\begin{equation*}\tag*{$[3]'$}B=4z-16z^3+56z^5-160z^7-\ldots\end{equation*}
\par\indent Gleichung [4] ergibt sich, indem man [3] in [1] einsetzt; demgemäß ist auch in [4] der Koeffizient von $z^6$ unrichtig; die richtige Formel lautet
\begin{equation*}\tag*{$[4]'$}1-4z^2+20z^4+663z^6+\ldots\end{equation*}
\par\indent Wenn wir die Bedeutung der Formeln [3], [4] in moderner Ausdrucksweise hervorheben wollen, so werden wir sagen: \textls[120]{diese Formeln stellen die Uniformisierung des vollständigen elliptischen Integrals zweiter Gattung als Funktion des Moduls} $B$ \textls[120]{mit Hilfe der Modulfunktion dar.} Freilich ist hier alles noch rein formal, da eine analytische Bestimmung der Veränderlichen
\Needspace{4\baselineskip}
\par\smallskip\noindent\rule{.25\linewidth}{.4pt}\par
{\small
\par\indent *) Solche Reihen betrachtet \textsc{Gauss} auch in der Tagebuchnotiz Nr. 58 vom 16. Febr. 1797; vergl. auch die Bemerkung \textsc{Scherings}, Werke III, S. 493, Zeile 4—6, wonach \textsc{Gauss} den Zusammenhang zwischen agM. und den Reihen, deren Exponenten die Quadratzahlen sind, 1794 gekannt haben soll. Mit dem agM. hat sich \textsc{Gauss}, wie er April 1816 an \textsc{Schumacher} schreibt (siehe oben Briefwechsel [7.], S. 247), seit 1791 beschäftigt. Wie \textsc{Gauss} in der Anzeige seiner Abhandlung \textit{Determinatio attractionis etc.} (Werke III, S. 360) angibt, hat er die Untersuchungen über das agM. unabhängig von den ähnlichen Untersuchungen von \textsc{Lagrange} (vergl. oben S. 251, Fußnote) und \textsc{Legendre} angestellt.
\par\indent **) Man vergl. z. B. die Tagebuchnotiz Nr. 49 vom 27. Dezember 1796 und die Methoden des Abschnittes [I.].
\par}
\clearpage
\fancyhead{}\fancyhead[R]{\small263}\fancyhead[C]{\scriptsize BEMERKUNGEN ZUM AGM. ABSCHNITT II.}
\SourceLeaf{270}
\noindent $z$ noch nicht gegeben wird. Diese ergibt sich erst später, siehe den Abschnitt [IV.]. Der Ausdruck [7] hat den Wert
\begin{equation*}\tag*{$[7]'$}\frac1{B^2}\left(\frac2\pi\frac qa-\frac{a}{M(a,b)}\right);\end{equation*}
\noindent aus der Gleichung (siehe \textit{Abriß}, Gleichung (33))
\[r^2=\frac{c}{M(a,b)}\]
\noindent folgt nämlich
\[\frac{a}{M(a,b)}=\frac{r^2}{B}.\]
\par\indent Abgesehen von dem Faktor $-\frac1{B^2}$ kommt der Ausdruck $[7]'$ auch in [19], oben S. 178, vor, ferner steht er mit der in [34], oben S. 180, auftretenden Größe $\mathbf L$ in der Beziehung
\begin{equation*}\tag*{$[7]''$}\frac{a}{M(a,b)}-\frac{2q}{\pi a}=\nu\mathbf L\qquad\left(\nu=B=\frac ca\right).\end{equation*}
\noindent Auf diese Beziehung kommen wir weiter unten (S. 265) zurück.
\par\indent Die Bedeutung der in den Formeln [11]—[16], [28], [29], [30], oben S. 178, 179, auftretenden Größe $\mathbf M$ ergibt sich am einfachsten wie folgt.
\par\indent In [35], oben S. 180, bedeutet $M$ das agM., und zwar ist
\[M\sqrt{1-\nu^2}=M(1,\sqrt{1-\nu^2})\,\text{\textsuperscript{*)}},\]
\noindent es ist also
\begin{equation*}\tag*{$[35]'$}K(\nu)=\frac1{M(1,\sqrt{1-\nu^2})}\,\text{\textsuperscript{**)}}.\end{equation*}
\par\indent Damit stimmt die Reihenentwicklung [33] (vergl. [2] und [6]), ferner die Funktionalgleichung [20], die in [32] und [39] wiederholt ist. Die Gleichung [38], worin $K'$ die Derivierte von $K$ nach $\nu$ bedeutet, kann zur Definition der Größe $L$ dienen. Aus der für $K'$ geltenden Funktionalgleichung [40] folgt dann in Verbindung mit [20] die in [29] angegebene Funktionalgleichung für $L$; unter derselben Nummer findet man die Definition von $\mathbf M$ als Quotienten von $L$ und $K$. Es ist also
\begin{equation*}\tag*{$[29]'$}\mathbf M(\nu)=1+\left(\nu-\frac1\nu\right)\frac{K'(\nu)}{K(\nu)}.\end{equation*}
\par\indent Mit der durch [38] definierten Größe $L$ steht die durch die Reihenentwicklung [34] erklärte Größe $\mathbf L$ in naher Beziehung\textsuperscript{***)}.
\par\smallskip\noindent\rule{.25\linewidth}{.4pt}\par
{\small
\par\indent *) Die Unterscheidung zwischen den Buchstaben $M$ und $\mathbf M$ ist in der \textsc{Gauss}schen Handschrift nicht gemacht.
\par\indent **) $K(\nu)$ unterscheidet sich also von dem \textsc{Jacobi}schen
\[K=\int_0^1\frac{dx}{\sqrt{(1-x^2)(1-\nu^2x^2)}}\]
\noindent nur durch den Faktor $\frac2\pi$.
\par\indent ***) Die Unterscheidung zwischen den Buchstaben $\mathbf L$ und $L$ ist auch vom Herausgeber interpoliert (vergl. Fußnote *)). Wahrscheinlich hat \textsc{Gauss} die Bezeichnung gegen eine »Uraufzeichnung« geändert. So hat z. B. die Gleichung [38] ursprünglich gelautet $\nu K+(\nu^2-1)K'=L$ und ist dann durch Korrekturen in die oben wiedergegebene Form gebracht worden.
\par}
\clearpage
\ArithmeticHead{264}{L}{ANALYSIS. NACHLASS.}
\SourceLeaf{271}
\par\indent Für $\mathbf L$ zeigt zunächst die Vergleichung von [34] mit der Reihenentwicklung (21) des \textit{Abrisses} für den Ellipsenquadranten $q$, daß
\begin{equation*}\tag*{$[34]'$}-\frac\pi2a\mathbf L=\frac{dq}{d\nu}\end{equation*}
\noindent ist (in (21) ist $k=\nu$ zu nehmen)\textsuperscript{*)}. Ferner führt die Vergleichung der Reihenentwicklungen [33], [34] und ihrer Derivierten nach $\nu$ zu den Differentialrelationen [36] oder [37]. Die Vergleichung der aus [37] folgenden Gleichung
\begin{equation*}\tag*{$[38]'$}\mathbf L=\nu K+(\nu^2-1)K'\end{equation*}
\noindent mit [38] gibt die Beziehung zwischen $\mathbf L$ und $L$:
\begin{equation*}\tag*{$[34]''$}L=\frac1\nu\mathbf L=\frac12+\frac12\,.\,\frac12\,.\,\frac34\nu^2+\frac{1\,.\,1\,.\,3\,.\,3\,.\,5}{2\,.\,2\,.\,4\,.\,4\,.\,6}\nu^4+\ldots,\end{equation*}
\noindent sodaß also (vergl. $[34]'$
\begin{equation*}\tag*{$[34]'''$}-\nu a\frac\pi2L=\frac{dq}{d\nu}\end{equation*}
\noindent ist\textsuperscript{**)}. — Für das durch $[29]'$ definierte $\mathbf M$ verifiziert man sofort die Funktionalgleichung [30] oder [12] ([11] ist einfach die Umkehrung von [12]); aus dieser läßt sich die Entwicklung [28] unmittelbar herleiten. Die Potenzreihenentwicklung [13]\textsuperscript{***)} ergibt sich entweder direkt durch Anwendung der Funktionalgleichung [12] oder durch Division von $[34]''$ mit [33]. Diese doppelte Herleitung von [13] bildet eine erwünschte Bestätigung unserer Deutung der \textsc{Gauss}schen Formeln.
\par\indent Wir bemerken noch, daß sich aus den Differentialrelationen [37] sofort die homogenen linearen Differentialgleichungen zweiter Ordnung herleiten ließen, denen $K$ und $\mathbf L$ als Funktionen von $\nu$ genügen. Diese Differentialgleichungen gehören (nach $[38]'$) zu derselben \textls[120]{Klasse} (im Sinne von \textsc{Riemann}); $K$ und $\mathbf L$ sind in der späteren Terminologie von \textsc{Gauss} (\textit{Disquisitiones circa seriem}, 1812, Werke III, S. 130) \textls[120]{functiones contiguae}.
\par\indent Nun kann die Darstellung [14] bezw. [15] des Ellipsenumfangs durch Vergleichung mit den Reihenentwicklungen [5] oder [17] verifiziert werden. Wir setzen eine Ableitung durch Integralformeln hierher.
\par\smallskip\noindent\rule{.25\linewidth}{.4pt}\par
{\small
\par\indent *) Dieselbe Entwicklung für $\frac{dq}{d\nu}$ findet sich auf Seite 129 der oben erwähnten \textsc{Euler}schen Abhandlung. \textsc{Euler} hat $a=1$. Die Bedeutung der Beziehung $[34]'$ wird weiter unten hervortreten.
\par\indent **) Wenn man die Integraldarstellung von $K$ beachtet, so erkennt man erst die wahre Bedeutung der von \textsc{Gauss} eingeführten Größe $L$. Aus $[34]'''$ folgt nämlich sofort
\[L=\frac2\pi\int_0^1\frac{x^2dx}{\sqrt{(1-x^2)(1-\nu^2x^2)}},\]
\noindent so daß $L$ gewissermaßen als ein Gegenstück zu
\[K=\frac2\pi\int_0^1\frac{dx}{\sqrt{(1-x^2)(1-\nu^2x^2)}},\]
\noindent erscheint. Es ist dies um so bemerkenswerter, als wie bereits erwähnt, \textsc{Gauss} zur Zeit, als er diese Aufzeichnungen machte, die Integraldarstellung von $K$ nicht gekannt hat.
\par\indent ***) In [13] ist, wahrscheinlich infolge eines Schreibfehlers, der Koeffizient von $x^6$ unrichtig; er lautet $\frac{41}{16048}$ (vergl. oben S. 188, art. [5.], Gleichung [28]).
\par}
\clearpage
\fancyhead{}\fancyhead[R]{\small265}\fancyhead[C]{\scriptsize BEMERKUNGEN ZUM AGM. ABSCHNITT II.}
\SourceLeaf{272}
\par\indent Aus den Integraldarstellungen von $q$ und $K$ folgen in bekannter Weise\textsuperscript{*)} die zuerst von \textsc{Legendre} (\textit{Exercices de calcul intégral} I, 1811, S. 61 u. ff.) aufgestellten Relationen (die sogenannten Klassenbeziehungen) zwischen den vollständigen Integralen erster und zweiter Gattung:
\begin{equation*}\tag*{$(\alpha)$}\frac\pi2\frac{dK}{d\nu}=\frac qa\frac1{\nu(1-\nu^2)}-\frac\pi2\frac K\nu,\end{equation*}
\begin{equation*}\tag*{$(\beta)$}\frac1a\frac{dq}{d\nu}=\frac1a\frac q\nu-\frac\pi2\frac K\nu,\end{equation*}
\noindent von denen die zweite $(\beta)$ vollständig mit der Gleichung $[7]''$ S. 263 übereinstimmt. Aus $(\alpha)$ ergibt sich
\[q=a\nu(1-\nu^2)\frac\pi2\frac{dK}{d\nu}+(1-\nu^2)\frac\pi2K,\]
\noindent und mit Rücksicht auf die Definition von $\mathbf M$ in $[29]'$:
\[q=\frac\pi2aK(1-\nu^2\mathbf M).\]
\noindent Diese Formel ist nichts anderes als die von \textsc{Gauss} aufgezeichnete Formel [14]. Wir haben sie noch einmal verifiziert, weil sich in den \textsc{Gauss}schen Formeln [14] und [15] eine Anzahl von Schreibfehlern befindet, die vielleicht daher rühren, daß \textsc{Gauss} in einer »Uraufzeichnung« $a$ statt $b$ geschrieben hatte (vergl. die Integralformel [10]). Die richtigen Gleichungen lauten für den Quadranten:
\begin{equation*}\tag*{$[14]'$}q=\frac\pi2aK\left(\sqrt{\frac{a^2-b^2}{a^2}}\right)\left\{1-\frac{a^2-b^2}{a^2}\mathbf M\left(\sqrt{\frac{a^2-b^2}{a^2}}\right)\right\},\end{equation*}
\begin{equation*}\tag*{$[15]'$}q=\frac\pi2a\frac{1-N\left(\sqrt{\frac{a^2-b^2}{a^2}}\right)}{\text{Med. inter }1\text{ et }\dfrac ba},\end{equation*}
\noindent wo $N$ durch die Gleichung [16] definiert wird\textsuperscript{**)}.
\par\indent Setzt man in $[14]'$ für $\mathbf M$ seinen Wert $\mathbf M=\frac LK$ oder, mit Rücksicht auf $[34]''$,
\[\mathbf M=-\frac2\pi\frac1{\nu a}\frac1K\frac{dq}{d\nu},\quad\nu=\sqrt{\frac{a^2-b^2}{a^2}},\]
\noindent ein, so erhält man für $q$ den Ausdruck
\begin{equation*}\tag*{$[14]''$}q=\frac\pi2aK(\nu)+\nu\frac{dq}{d\nu},\end{equation*}
\noindent was mit der zweiten der \textsc{Legendre}schen Gleichungen $(\beta)$, also mit $[7]''$ übereinstimmt.
\par\smallskip\noindent\rule{.25\linewidth}{.4pt}\par
{\small
\par\indent *) Vergl. z. B. H. \textsc{Durège}, \textit{Theorie der elliptischen Funktionen} (Leipzig 1878), S. 290.
\par\indent **) Während uns heute die Beziehung des agM. zum vollständigen Integral erster Gattung als die am nächsten liegende erscheint, tritt bei den älteren Analysten (\textsc{Landen}, \textsc{Lagrange}) die Beziehung zum vollständigen Integral zweiter Gattung immer zuerst hervor. Auch \textsc{Schumacher} hat, als er sich gelegentlich (Brief vom 5. April 1816, siehe oben Briefwechsel [6.], S. 246) mit dem agM. beschäftigte, als erstes Ergebnis die Formel $[15]'$ gefunden. Den Zusammenhang des agM. mit dem Ellipsenquadranten hat \textsc{Gauss} in der »\textit{Determinatio attractionis}« (1818) (Werke III, S. 354 oben) und in der Anzeige dieser Abhandlung (ebenda, S. 360) nach andern Methoden abgeleitet. Vergl. auch die aus 1800 stammende nachgelassene Abhandlung über das agM. (Werke III, S. 261 ff., besonders S. 373 die Formel für $dM(x,y)$, die auch in der Aufzeichnung [VI.], oben S. 208 vorkommt).
\par}
\noindent X1.\hfill34\par
\clearpage
\ArithmeticHead{266}{L}{ANALYSIS. NACHLASS.}
\SourceLeaf{273a}
\par\indent Die Aufzeichnung [6.] S. 180 steht mit den in den artt. [1.]—[5.] enthaltenen in keinem unmittelbaren Zusammenhang; sie ist darum bemerkenswert, weil sie zeigt, daß \textsc{Gauss} schon damals neben dem einfachen Algorithmus aus $a,b$ auch den aus $a$ und $\sqrt{a^2-b^2}$ sowie die aus beiden durch Rückwärtsverlängerung hervorgehenden Ketten (vergl. \textit{Abriß}, Artikel 1. S. 252) betrachtet hat. Das Zahlenbeispiel hängt mit $M(\surd2,1)$ (vergl. die Tagebuchnotiz Nr. 98, oben S. 260) zusammen.
\par\bigskip
\par\indent Die in dem Abschnitt [II.] zusammengestellten Leisteaufzeichnungen beginnen bei Seite 25. Da die den Tagebuchnotizen von 1796 entsprechenden Leistenotizen nur etwa bis zu der Seite 16 reichen, so ergibt sich also als untere Zeitgrenze für alle hier wiedergegebenen Aufzeichnungen das Jahr 1797; eine genaue Zeitbestimmung werden wir erst nach Besprechung der Abschnitte [III.] und [IV.] geben können. Die Notizen [II.] gehören zu den »irregulären« Eintragungen (vergl. oben S. 170, Fußnote *)), da sie mit ersichtlich späterer Schrift und Tinte auf zerstreuten Stellen, die bei den regulären Aufzeichnungen unbeschrieben geblieben waren, eingetragen sind. Ihre Form ist durchaus fragmentarisch, sie zeigen zahlreiche Schreib- und Rechenfehler (vergleiche z. B. [13], [14], [3], [4]) und Folgewidrigkeit in der angewandten Bezeichnung ($M,L$ in verschiedenen Bedeutungen, vergl. oben S. 263). Wir haben also in diesen Notizen nicht eine redigierte Darstellung fertiger Ergebnisse vor uns, sondern Rechnungen, die während der Arbeit, vielleicht zum Teil im Anschluß an ältere, für uns wohl verlorene Notizen gemacht worden sind. Daß solche ältere Notizen vorhanden waren, ist unzweifelhaft, da ja (vergl. oben S. 262 Fußnote *)) \textsc{Gauss} sich seit 1791 mit dem agM. beschäftigt hatte; auch hat die Angabe \textsc{Scherings}, daß \textsc{Gauss} schon 1794 den Zusammenhang des agM. mit den Reihen, deren Exponenten die Quadratzahlen sind, gekannt habe, sehr große innere Wahrscheinlichkeit (vergl. die Tagebuchaufzeichnungen Nr. 58 vom 16. Februar 1797 und Nr. 7 vom 24. Mai 1796). Ihrem Inhalte nach enthalten die Aufzeichnungen [II.] wesentlich drei Elemente:
\par\indent 1) Die Beziehung des agM. zum Ellipsenquadranten,
\par\indent 2) die Beziehung des agM. zu den Reihen, deren Exponenten die Quadratzahlen sind und die dadurch bedingte Einführung der Größe $z$,
\par\indent 3) die grundsätzliche Einführung des reziproken Wertes $K$ des agM. (vergl. oben S. 260) und seine Reihenentwicklung nach positiven ganzen Potenzen von $\nu$ (siehe [2], [6], [18], [33]), die den ganzen Notizen gleichsam als Unterbau zu Grunde liegt.
\clearpage
\endgroup

